% Khoi TIENG VIET (bat buoc theo muc 3 quy dinh HUJOS-TT:
% "Authors should provide an abstract and keywords in Vietnamese and English").
%
% BAI DANG NOP BANG TIENG ANH, nen dat ban tieng Anh o DAU bai (dung nhu bai
% HUJOS-TT da xuat ban that, ma 7999, DOI 10.26459/hueunijtt.v135i2B.7999, co
% citation_language = en va 0 doan tieng Viet tren trang) va ban tieng Viet o
% CUOI bai. Quy dinh van ban viet theo chieu nguoc lai (Viet o dau, Anh o cuoi)
% vi no mac dinh bai viet bang tieng Viet; day la anh xa doi xung, giu du ca hai
% ngon ngu theo dung yeu cau "Vietnamese and English".
%
% Tu khoa: 3-5 cum, chi dung THUA TU CO TRONG TITLE cua nguon da verify
% (khong dung khai niem tu dat). Kiem bang grep tren refs.bib:
%   "alternate wetting and drying"  -> 5 title
%   "MRV" / "monitoring reporting verification" -> 4 title
%   "carbon credit"                 -> 4 title
%   "low-emission rice"             -> 2 title
%   "digital twin"                  -> 2 title
% Loai bo "fraud detection", "adversarial", "integrity": 0 title nao co.

\section*{Tóm tắt và từ khóa (Tiếng Việt)}

Đo đạc, báo cáo, thẩm định (MRV) là nút thắt khiến phần giảm phát thải của lúa phát thải thấp chưa thành tín chỉ các-bon, vì giám sát chiếm khoảng ba phần tư chi phí. Trong tưới khô ướt luân phiên (AWD), mỗi hành động đủ điều kiện cấp tín chỉ là một lệnh bật máy bơm, nên nhật ký nông dân tự ghi trông như bằng chứng miễn phí, xếp ở bậc ưu tiên cao nhất theo Quyết định 4801/QĐ-BNNMT. Bài báo dựng lớp kiểm toán đó trên bản sao số (digital twin) FAO-56, phát lại nhật ký với ERA5 công khai rồi cho đối thủ tối ưu hoá tấn công. Kết quả, mọi bộ phát hiện dựa trên quỹ đạo phát lại có điểm mù tất yếu, kèm cơ chế khắc phục có chi phí biên gần bằng không. Tầng vật lý phát hiện khai báo vượt sức chứa ẩm đồng ruộng, nhưng bỏ lọt đối thủ dịch thời điểm lệnh, giữ nguyên tổng lượng, khiến số pha khô được cấp tín chỉ tăng $33\%$ (ghép cặp). Bộ phát hiện mô hình điểm kết tội oan $26{,}7\%$ nông dân trung thực; bộ bền vững hạ oan còn $5{,}0\%$ nhưng lọt phần lớn tấn công dịch thời điểm. Trên $281$ nhật ký ở ba trạm Đồng bằng sông Cửu Long, phép neo bắt $116/116$ ca giả mạo, không oan ai. Vật lý do đó chỉ ràng buộc nhật ký; dấu thời gian phải do thiết bị cấp.

\textbf{Từ khóa:} tưới khô ướt luân phiên (AWD), đo đạc báo cáo thẩm định
(MRV), tín chỉ các-bon, lúa phát thải thấp, bản sao số.

\vspace{4pt}

{\raggedright\normalfont\fontsize{9}{11}\selectfont
\textbf{Thông tin tác giả.} Mai Xuân Văn$^{1}$, Nguyễn Tương Tri$^{2,3,*}$\\[1pt]
$^{1}$ Trường Trung học phổ thông (THPT) Thuận Hóa, Trường Đại học Sư phạm, Đại học Huế;
$^{2}$ Viện Đào tạo Mở và Công nghệ thông tin, Đại học Huế;
$^{3}$ Trường Đại học Sư phạm, Đại học Huế, 34 Lê Lợi, Huế, Việt Nam\\[1pt]
Học hàm, học vị: (điền); điện thoại: (điền); E-mail: ntuongtri@hueuni.edu.vn;
$^{*}$ Tác giả liên hệ: Nguyễn Tương Tri.\\[2pt]
Nghiên cứu này được tài trợ bởi Đại học Huế trong đề tài mã số DHH2025-19-07.
\hfill\textit{Chữ ký tác giả:} \rule{3.5cm}{0.4pt}\hfill\rule{3.5cm}{0.4pt}\hspace*{\fill}
\par}
